% ch1-introduction.tex -- part of thesis_v3.tex; do not add \documentclass here.
% Draft with the thesis-chapter workflow. Unsupported claims must use \fillin{...}.

\chapter{Introduction}
\label{ch:introduction}

\section{Motivation and Problem Statement}
\label{sec:motivation_problem}

A printed circuit board's power delivery network (PDN) has to supply the integrated circuits on
the board with a stable voltage while they draw current that changes quickly. The usual way to
express this requirement in the frequency domain is the target impedance rule: the impedance
$Z(f)$ seen looking into the PDN from the load must stay below a frequency-dependent limit
$Z_{\mathrm{target}}(f)$ across the band of interest \cite{smith1999target}. Decoupling
capacitors (decaps) are the main tool the designer has for shaping $Z(f)$, and the design task is
to decide how many of them to place and where.

In this thesis the band of interest is $1$--$600$\,MHz, sampled at $231$ discrete frequency bins.
The bin locations are fixed by the file \texttt{configs/Frequency\_data\_hz.npy}, which holds
$231$ values running from $1000000.0$\,Hz to $600000000.0$\,Hz. The grid is not uniformly spaced:
the spacing between successive bins takes the three values $0.1$\,MHz, $1$\,MHz and $10$\,MHz.
\fillin{state why the frequency grid uses three different step sizes and who chose the
breakpoints -- the ECADSTAR sweep setup or export script that produced Frequency\_data\_hz.npy
would settle it}

The impedance limit itself is stored as an array of shape $(231,1)$ in
\texttt{configs/target\_impedance.npy}, so there is one limit per bin. It is not a constant: the
first element is $0.8$ and the last is $50$ (stored as $49.99999999999999$ in \texttt{float64}),
and the values vary with frequency in between. Stage~2 of the framework loads this file, casts it
to \texttt{float32} and squeezes it to a length-$231$ vector before use.
\fillin{confirm the unit of configs/target\_impedance.npy. The rest of this thesis reads the mask
as ohms because the Stage-2 code compares it against a tensor named \texttt{imp\_ohm} and its
default selection metric is named \texttt{max\_ohm}, but that is an inference from variable
naming, not a unit recorded with the array; a design note or the ECADSTAR export settings would
settle it}
\fillin{record the provenance of configs/target\_impedance.npy -- who specified this mask, and
from which load-current or ripple budget it was derived}

\subsection*{The placement problem}

The board studied here exposes a fixed set of $52$ candidate decap sites. These are the valid
cells of a $7\times 8$ physical grid: the code sets \texttt{OCC\_GRID\_H}, \texttt{OCC\_GRID\_W}
$= 7, 8$ and marks four cells invalid, namely $(0,3)$, $(0,4)$, $(6,3)$ and $(6,4)$, leaving $52$
slots labelled \texttt{C1}--\texttt{C52}. A design decision is therefore a binary placement
vector
%
\begin{equation}
  \mathbf{b} \in \{0,1\}^{52},
  \qquad
  K = \lVert \mathbf{b} \rVert_0 ,
  \label{eq:placement_vector}
\end{equation}
%
where $b_i = 1$ means a decap is mounted in slot $i$ and the budget $K$ is the number of decaps
used. $K$ matters because decaps cost money and board area, so a placement that meets the mask
with fewer parts is preferable to one that needs more.

Two properties make this problem hard. First, the map $\mathbf{b} \mapsto Z(\cdot)$ is only
available through an expensive field solver, so every candidate placement that is to be judged
exactly has to be simulated. Second, the discrete space has $2^{52}$ elements, which is stated in
the project documentation as intractable to enumerate. The combination means that neither
exhaustive search nor a simulation-in-the-loop search that evaluates a large number of candidates
is practical.

\subsection*{The conventional design loop}

The conventional workflow this thesis targets is iterative. The engineer places a set of decaps,
simulates the resulting impedance spectrum, and checks it against the mask. If peaks violate the
mask, a spatial power-integrity distribution is computed at the offending frequency, decaps are
moved towards the hotspots that this distribution reveals, and the cycle repeats. Every turn of
that loop contains at least one full solver run.

\fillin{quantify the cost of this loop for the board under study -- engineer hours per iteration
and the number of iterations typically needed to converge; nothing in the repository records
this, and the literature search found no freestanding paper that measures it}
The claim that iterative and manual decap tuning is slow enough to motivate a learned replacement
is asserted in the motivation sections of the machine-learning and genetic-algorithm decap
literature rather than measured in a dedicated study
\cite{swaminathan2020demystifying,zhang2020enhanced,ars2023ai,convga2024gcn}.

Ground-truth impedance data for this thesis is produced with ECADSTAR, which runs on a Windows
host addressed from WSL; the active-learning configuration carries absolute Windows paths for the
ERF, EMC and PEB locations, for example
\texttt{C:\textbackslash Users\textbackslash muthusamy\textbackslash Desktop\textbackslash
design\textbackslash H-shape.emc\textbackslash H-shape.erf}. The configuration allocates a wait
timeout of $18000$ seconds for a simulation batch and $900$ seconds for batch start, polled every
$40$ seconds. These are configured timeouts and not measurements.
\fillin{measure and report the wall-clock time of one ECADSTAR PI simulation, per layout and per
batch; a timing log from an AL cycle's simulate step would settle it}
\fillin{report surrogate inference time on the same hardware and the resulting speed-up factor
over simulation -- no artifact in the repository contains this comparison}

\subsection*{Learned approaches to decap placement}

Learning-based decap placement is not new. Deep reinforcement learning was applied to decap
selection from 2019 onwards \cite{zhang2019decoupling}, extended with an impedance-matrix
reduction method and double Q-learning to cut optimisation time \cite{zhang2020enhanced}, and
applied to silicon-interposer 2.5-D/3-D IC decap design in the same period
\cite{park2020drldecap}. The same line later moved to transformer-based policies
\cite{park2022transformer,kim2023devformer}, and separate strands combine graph neural networks
with genetic algorithms \cite{convga2024gcn} or use learned models purely as fast impedance
predictors \cite{zhang2021fast}. Chapter~\ref{ch:theory_background} treats this work in detail
and states where the present framework differs.

\subsection*{Approach taken in this thesis}

The framework documented in this repository replaces the conventional loop with four elements.
A learned differentiable surrogate of one fixed board, built as a multi-input variational
autoencoder, stands in for the field solver. Gradient-based inverse design in the latent space of
that surrogate searches for placements under a budget $K$. An active-learning loop optionally
refines the surrogate on layouts where its error or uncertainty is high, simulating only a
selected batch rather than every candidate. Finally, a performance-analysis layer decodes a
frequency-conditioned spatial heatmap at the peak frequencies of a predicted spectrum. The
project documentation gives the purpose of that layer as hotspot interpretation, that is, the
intent is to inspect a predicted violation spatially instead of only in the spectrum. Whether the
decoded heatmaps are accurate enough for that use is evaluated in
Section~\ref{subsec:val_pi_heatmaps}; nothing in this chapter asserts that they are.

The surrogate models three modalities jointly for a single board: decap occupancy as the binary
vector of Equation~\eqref{eq:placement_vector}, the PI spectrum magnitude over the $231$
frequency bins, and a $64\times 64$ spatial PI heatmap at selected MHz anchors. Conditioning
inputs are the budget $K$ and the inspection frequency. In the current model track,
\texttt{exp059\_capacity\_freq}, the three modalities carry the shapes $(B,1,64,64)$ for the
heatmap, $(B,52)$ for occupancy and $(B,1,231)$ for the impedance vector, where $B$ is the batch
size. These shapes are the ones documented in the docstring of the model's \texttt{encode}
method in \texttt{experiments/exp059\_capacity\_freq/codes/vae\_multi\_input\_simple.py}, which
lists the encoder's three modality inputs. $K$ is defined in code as the number of ones in the
occupancy vector, so $K \in \{0,\dots,52\}$.
\fillin{state the exact return signature of the exp059 \texttt{forward} pass and the form of its
occupancy output -- in particular whether that output is decoder logits that are thresholded
before use, or an already-binarised vector. The shapes above come from the \texttt{encode}
docstring and are input shapes; the \texttt{forward} and \texttt{decode} definitions in the same
file would settle what is returned}

The inverse-design objective is stated per budget. For each $K$, a placement is feasible if the
predicted spectrum stays below the mask with a safety margin $m$,
%
\begin{equation}
  \hat{Z}(f; \mathbf{b}) \;\le\; Z_{\mathrm{target}}(f) - m
  \qquad \text{for all } f \text{ in the } 231 \text{ bins},
  \label{eq:feasibility_intro}
\end{equation}
%
where $\hat{Z}$ denotes the surrogate's predicted impedance as opposed to the simulated $Z$.
Among feasible candidates the one with the lowest peak impedance $\max_f \hat{Z}(f;\mathbf{b})$
is selected; the implementation's default selection metric is \texttt{max\_ohm}, the maximum of
\texttt{imp\_ohm} over the frequency bins. If no candidate for a given $K$ is feasible, the
procedure records an explicit \texttt{no\_solution.json} for that $K$ instead of returning a
best-effort placement. As the code stands, the Stage-2 script
\texttt{pipelines/latent/optimize.py} loads its occupancy decoder and impedance surrogate from
the legacy experiment \texttt{exp038\_true\_multi} rather than from
\texttt{exp059\_capacity\_freq}, so the objective above is stated as method only and no
inverse-design outcome is attributed to the current surrogate anywhere in this thesis.

One implementation detail of the search is load-bearing and is introduced here because it shapes
the whole formulation. Reading a placement out of the decoder requires a hard top-$K$ selection,
which is a discrete operation with no useful gradient. The straight-through operator used in
Stage~2 returns the hard top-$K$ binary vector in the forward pass but passes the gradient of the
soft occupancy probabilities backwards; its docstring states that without this, the top-$K$
argsort severs the gradient and the latent vector never moves in response to the spectrum target.

\subsection*{Scope and validity boundary}

The scope is deliberately narrow. The surrogate is trained for a single board and one design
family, with no cross-board generalisation claimed; the decap type and the $52$-slot set are
fixed, so no multi-value or multi-type catalogue is covered; and the current training filter
keeps layouts with $K \le 30$. The normalised training dataset records $24979$ unique layouts and
$582997$ manifest rows.

Surrogate spectra and heatmaps are predictions, not sign-off. A proposed design has to be
re-simulated with the ground-truth PI flow before it can be accepted. This boundary is stated
explicitly in the project's own validity notes and is repeated here because it constrains every
result claim made later in the thesis.

\fillin{describe the board under study -- layer stack, physical dimensions, IC port location and
the power bus being decoupled; the AL configuration references an ERF named H-shape and a
powerbus named Power\_GND, but no board specification document exists in the repository}
\fillin{state the capacitance value and part type assumed for the 52 slots in exp059; the only
decap-profile artifact found belongs to a legacy experiment note and not to the exp059 track}
\fillin{personal and industrial motivation for choosing this topic -- not recorded anywhere in
the repository, must be written by hand}

\figplace{Board slot map: the $7\times 8$ grid with the four invalid cells removed, slots labelled C1--C52.}

\figplace{Conventional place--simulate--inspect--adjust loop next to the proposed surrogate-based loop.}

\section{Research Objectives and Contributions}
\label{sec:objectives_contributions}

The overall aim follows from Section~\ref{sec:motivation_problem}: to replace the
simulate-in-the-loop decap placement search on one fixed board with a learned differentiable
surrogate, to perform the search by gradient descent in that surrogate's latent space under a
budget constraint, and to add an analysis layer that localises impedance violations spatially.
Active learning is included as a mechanism for spending a limited simulation budget on the
layouts where the surrogate is least reliable.

Beyond that aim, the research questions and the contribution list are not fixed yet. No research
questions and no enumerated contributions exist in the repository, in the documentation, or in
the thesis vault, and they require supervisor agreement before they are written down.

\fillin{RQ1--RQn: formulate the research questions with the supervisor. Candidate themes visible
in the artifacts are (i) whether a multi-modal VAE surrogate of one board is accurate enough to
drive placement decisions, (ii) whether latent-space gradient search through a straight-through
top-$K$ read-out finds feasible placements at smaller $K$ than the alternatives, and (iii)
whether active learning improves the surrogate where it matters under a fixed simulation budget.
These are candidates only and must not be treated as agreed.}

\fillin{Contribution list: enumerate the claimed contributions once the experiments in
Chapter~\ref{ch:results_evaluations} are complete and agreed with the supervisor. This list must
match Section~\ref{sec:summary_contributions} in the conclusion. Nothing in the repository
currently enumerates contributions.}

\fillin{State explicitly which of the objectives are demonstrated end-to-end and which are not.
As the code stands, the Stage-2 inverse-design script loads occupancy-decoder and
impedance-surrogate checkpoints from the legacy experiment exp038\_true\_multi, while the current
surrogate track is exp059\_capacity\_freq -- so no end-to-end inverse-design result may be
attributed to exp059 until Stage 2 is rewired. Resolve before claiming any objective as met.}

\fillin{Confirm the delimitation sentence: what this thesis does not attempt (signal integrity
sign-off, EM/EMC field solving, emissions, cross-board transfer). Wording to be agreed.}

\section{Structure of the Thesis}
\label{sec:thesis_structure}

The thesis has five chapters. Front matter provides a table of contents, a list of figures, a
list of tables, an unnumbered chapter on mathematical notation and an unnumbered list of
acronyms.

Chapter~\ref{ch:theory_background}, \emph{Theoretical Background}, sets out the theory the rest
of the work rests on. Section~\ref{sec:pdn_optimization} covers power delivery networks and the
optimisation problem, with subsections on impedance and decoupling
(\ref{subsec:impedance_decoupling}), PI simulation and its multimodal outputs
(\ref{subsec:pi_simulation_multimodal}), and the formal decap placement formulation
(\ref{subsec:decap_placement_formulation}). Section~\ref{sec:generative_latent_opt} covers
generative modelling and latent-space optimisation, with subsections on multimodal variational
autoencoders (\ref{subsec:vae_multimodal}), physics surrogates (\ref{subsec:physics_surrogate})
and latent-space optimisation (\ref{subsec:latent_space_opt}).
Section~\ref{sec:related_work} reviews related work on PDN optimisation methods
(\ref{subsec:pdn_opt_methods}) and generative learning approaches
(\ref{subsec:generative_learning_approaches}) and closes with the research gap
(\ref{subsec:research_gap}).

Chapter~\ref{ch:design_implementation}, \emph{Design and Implementation}, describes what was
built. Section~\ref{sec:data_prep_reprocessing} covers data preparation and reprocessing: the raw
PI outputs (\ref{subsec:raw_pi_outputs}), the multi-frequency dataset
(\ref{subsec:multifreq_dataset}) and the splitting and sampling strategy
(\ref{subsec:data_splitting_sampling}). Section~\ref{sec:multimodal_vae_surrogate} describes the
multi-modal variational autoencoder used as the generative surrogate, covering the architectural
design (\ref{subsec:architectural_design}), the training objective and physics regularisation
(\ref{subsec:objective_physics_reg}), training and inference (\ref{subsec:training_inference})
and the auxiliary guidance network (\ref{subsec:aux_guidance_network}).
Section~\ref{sec:latent_opt_decap} formulates the latent-space optimisation of decap placement
(\ref{subsec:opt_problem_formulation}) and the optimisation procedure with its feasibility test
(\ref{subsec:opt_procedure_feasibility}). Section~\ref{sec:experimental_setup} states the
experimental setup and Section~\ref{sec:active_learning} the active-learning loop used to scale
the model.
\fillin{state the per-cycle active-learning budget in Chapter 3 and resolve the conflict between
active\_learning\_pi/config/exp059\_gp\_error.json (candidates\_per\_k 1600, worst\_per\_k 80,
total\_ecad\_per\_cycle 640, K = 3..10) and docs/active-learning.md:110--112, which documents
8 x 400 candidates and 8 x 10 = 80 ECAD simulations for the older exp057 loop}

Chapter~\ref{ch:results_evaluations}, \emph{Results and Evaluations}, reports the measurements.
Section~\ref{sec:surrogate_performance} evaluates the generative surrogate through reconstruction
performance (\ref{subsec:reconstruction_perf}), validation of predicted PI heatmaps
(\ref{subsec:val_pi_heatmaps}) and evaluation of the auxiliary guidance network
(\ref{subsec:eval_aux_guidance}). Section~\ref{sec:latent_design_opt_results} reports the
latent-space design optimisation results, covering scalability and efficiency
(\ref{subsec:eval_scalability_efficiency}) and the spatial and capacitor-type evaluation
(\ref{subsec:spatial_cap_type_eval}). Section~\ref{sec:limitations} states the limitations and
Section~\ref{sec:ablations} the ablations.

Chapter~\ref{ch:conclusion_outlook}, \emph{Conclusion and Outlook}, summarises the contributions
(\ref{sec:summary_contributions}) and gives directions for future research
(\ref{sec:future_research}).

Three appendices follow the bibliography. Appendix~\ref{ch:appendix_simulation_dataset}
documents the power integrity simulation environment and the dataset specification.
Appendix~\ref{ch:appendix_hyperparameters} lists the hyperparameter configuration used for
training and optimisation. Appendix~\ref{ch:appendix_supp_results} collects supplementary
experimental results.
\fillin{confirm what Appendix C will contain; it is still declared inline in thesis\_v3.tex with
the note that it depends on experiments not yet run}

% ---------------------------------------------------------------
% FILL IN checklist:
%   1. Frequency-grid design rationale (three step sizes 0.1/1/10 MHz) -- resolved by the ECADSTAR
%      sweep setup or the export script that produced configs/Frequency_data_hz.npy.
%   2. Provenance of configs/target_impedance.npy (who specified the mask, from which current or
%      ripple budget) -- resolved by a design note or supervisor statement.
%   2b. Unit of configs/target_impedance.npy -- the ohms reading is an inference from the names
%      imp_ohm / SELECT_METRIC = "max_ohm" in pipelines/latent/optimize.py, not a unit stored with
%      the array. Now marked in the body as a \fillin rather than asserted.
%   3. Cost of the conventional place-simulate-inspect-adjust loop (engineer hours, iteration
%      count) -- no repository artifact and no freestanding paper measures it; needs either an
%      industrial source or removal of the quantitative framing.
%   4. Measured ECADSTAR wall-clock per layout and per batch -- resolved by a timing log from an
%      AL simulate step. Only configured timeouts (18000 s / 900 s / 40 s poll) exist today.
%   5. Surrogate inference time and speed-up factor vs simulation -- no artifact located; needs a
%      benchmark run on stated hardware.
%   6. Board specification (layer stack, dimensions, IC port, power bus) -- only "H-shape.erf" and
%      powerbus "Power_GND" appear, in the AL config; needs the ECADSTAR project or a board note.
%   7. Decap value/part type for the 52 slots in exp059 -- configs/decap_profiles.npy is referenced
%      only from a legacy experiment note (exp031), not the exp059 track.
%   8. Personal/industrial motivation -- not recorded anywhere; hand-written by Kovarthanam.
%   9. RQ1--RQn -- require supervisor agreement; no research questions exist in docs/,
%      thesis_vault/00-Thesis/ or the chapter files. Candidate themes listed in the marker are
%      suggestions only.
%  10. Contribution list -- no artifact enumerates contributions; ch5-conclusion.tex
%      sec:summary_contributions is empty. Must match once written.
%  11. Which objectives are demonstrated end-to-end -- blocked by the exp038/exp059 Stage-2 wiring
%      issue (item 2 of the evidence warnings).
%  12. Delimitation sentence wording (what the thesis does not attempt).
%  13. Per-cycle AL budget in Ch.3 -- conflict between exp059_gp_error.json and
%      docs/active-learning.md:110-112 named in the text.
%  14. Appendix C contents -- depends on experiments not yet run.
%  15. exp059 forward-pass return signature and the form of the occupancy output (logits vs.
%      binarised). The (B,1,64,64)/(B,52)/(B,1,231) shapes in the body are taken from the encode()
%      docstring, i.e. encoder INPUT shapes; forward() and decode() in
%      experiments/exp059_capacity_freq/codes/vae_multi_input_simple.py would settle what is
%      returned. Sourcing corrected in this revision; the numbers themselves were not changed.
%
% Citation verification queue:
%   - smith1999target : used here for the target-impedance rule. A prior audit found the stored DOI
%     pointed at a different 1999 paper; the refs.bib entry was repaired in a separate pass to DOI
%     10.1109/6040.784476, pages 284--291, checked against Crossref. The citation key is unchanged.
%     Re-read refs.bib before submission to confirm the repair is present and the VERIFY note there
%     reflects the current state. A number of other refs.bib entries still carry the generic
%     "VERIFY: bibliographic fields not confirmed from source page" note; count them against the
%     current file rather than trusting any figure quoted here.
%   - zhang2019decoupling : literature block records verified = "partial"; IEEE Xplore returned 403,
%     title/authors/venue confirmed from a co-author's publication list, page numbers unconfirmed.
%   - zhang2020enhanced, park2020drldecap : verified = "yes" in the literature block.
%   - swaminathan2020demystifying, ars2023ai, convga2024gcn, park2022transformer, kim2023devformer,
%     zhang2021fast : taken from the literature block's "known" list; bibliographic fields not
%     re-verified in this pass.
%   - The cited group {swaminathan2020demystifying, zhang2020enhanced, ars2023ai, convga2024gcn}
%     supports only the assertion that iterative decap tuning is slow, made in those papers'
%     motivation sections. It is NOT a measured benchmark and must not be paraphrased as one.
%   - Unit of configs/target_impedance.npy: no longer asserted in the body. See FILL IN item 2b.
%
% Evidence warnings:
%   1. Current model track: docs/cursor-handoff.md:6,:19 and docs/README.md:7-8 name
%      exp059_capacity_freq as current, while docs/framework-overview.md:40 and
%      docs/model-architecture.md:5 still name exp057_structured_graph. The chapter follows
%      exp059 per the constitution's project traps; the exp057 statements are stale. Root
%      README.md was not used as a source.
%   2. Stage 2 (pipelines/latent/optimize.py:45-47) sets EXPERIMENT = "exp038_true_multi" and loads
%      that experiment's checkpoints. No end-to-end inverse-design result may be attributed to
%      exp059. This chapter therefore describes the Stage-2 objective and the straight-through
%      top-K operator as method only, never as a demonstrated outcome. Since the audit this is
%      also stated in the RENDERED body, at the end of the Stage-2 objective paragraph in
%      Section 1.1, not only in Section 1.2's \fillin and in this comment.
%   2c. The performance-analysis (heatmap) layer is described as intent taken from
%      docs/framework-overview.md:25-30, not as a working capability; heatmap accuracy is deferred
%      to subsec:val_pi_heatmaps, which is currently a heading only.
%   3. exp059 capacity-phase modality probabilities conflict: docs/cursor-handoff.md:37 says
%      layout_train_prob = 0 and occ_only_encode_prob = 0; config.yaml holds 0.65 and 0.4.
%      Deliberately not used in Chapter 1; must be resolved before Chapter 3 quotes either.
%   4. AL per-cycle budget conflict (config vs docs/active-learning.md) -- named in the text as
%      a \fillin rather than resolved.
%   5. Chapter count conflict: thesis_vault/00-Thesis/Thesis Outline.md organises the material into
%      at least 7 chapters; Thesis_report/thesis_v3.tex has 5 and is authoritative per the source
%      hierarchy. Section 1.3 follows thesis_v3.tex.
%   6. Dataset anomaly not mentioned in this chapter but relevant later: dataset_meta.json records
%      samples_per_mhz = 8480 for the 80 MHz anchor while all other anchors have 24979.
%   7. No experimental outcome is stated anywhere in this chapter, by design -- it is written
%      before the remaining experiments.
%   8. Parameter count for exp059 is unknown; the ~10.7M figure in docs/model-architecture.md:5
%      belongs to exp057 and was deliberately not carried over.
% ---------------------------------------------------------------
